\documentclass[conference]{IEEEtran}

\usepackage{fontspec}
\setmainfont{Times New Roman}

\usepackage{amsmath,amssymb,amsfonts}
\usepackage{graphicx}
\usepackage{textcomp}
\usepackage{xcolor}
\usepackage{booktabs}
\usepackage{enumitem}
\usepackage{xurl}
\usepackage[hidelinks]{hyperref}
\usepackage{flushend}
\usepackage{sc26repro}

% IEEE conference \subsection* is italic at body size and blends into the
% text. Keep numbered I/II/A/B heads in IEEE style; make the AD template's
% unnumbered workflow heads bold with a small gap above.
\newcommand{\ADhead}[1]{%
  \par\vspace{0.75ex}%
  \noindent{\normalfont\normalsize\bfseries #1}\par
  \nobreak\vspace{0.3ex}\noindent\ignorespaces}
\newcommand{\ADpart}[1]{%
  \par\vspace{0.5ex}%
  \noindent{\normalfont\normalsize\bfseries #1}\par
  \nobreak\vspace{0.2ex}\noindent\ignorespaces}
\renewcommand{\artrel}{\ADhead{Relation To Contributions}}
\renewcommand{\artexp}{\ADhead{Expected Results}}
\renewcommand{\arttime}{\ADhead{Expected Reproduction Time (in Minutes)}}
\renewcommand{\artin}{\ADhead{Artifact Setup (incl. Inputs)}}
\renewcommand{\artinpart}[1]{\ADpart{#1}}
\renewcommand{\artcomp}{\ADhead{Artifact Execution}}
\renewcommand{\artout}{\ADhead{Artifact Analysis (incl. Outputs)}}

\setlist{nosep,leftmargin=1.15em,itemsep=2pt}
\setlist[description]{nosep,itemsep=3pt,leftmargin=1.7em,font=\normalfont\bfseries}

% Final submission: do NOT enable explanation/example tags.
% \usetag{explanation}
% \usetag{example}

\begin{document}

\twocolumn[%
{\begin{center}
\Huge
Appendix: Artifact Description\\
{\normalsize PDSW'26 --- WarpX / ADIOS2 I/O benchmarking study}
\end{center}}
]

%%%%%%%%%%%%%%%%%%%%%%%%%%%%%%%%%%%%%%%%%%%%%%%%%%%%%
%  AD Appendix only (AE deferred until after acceptance)
%%%%%%%%%%%%%%%%%%%%%%%%%%%%%%%%%%%%%%%%%%%%%%%%%%%%%

\appendixAD

\section{Overview of Contributions and Artifacts}

\subsection{Paper's Main Contributions}

\begin{description}
\item[$C_1$] Balanced-workload write-path characterization on Frontier and
  Perlmutter (weak scaling; multi-TB/s, including Blosc2).
\item[$C_2$] ADIOS2 aggregation (\texttt{TwoLevelShm}, \texttt{EveryoneWritesSerial},
  \texttt{DataSizeBased}) under extreme writer-side skew.
\item[$C_3$] \texttt{FlattenSteps} for lab-frame BTD snapshots from non-contiguous
  boosted-frame data at negligible write-time overhead.
\item[$C_4$] \texttt{SPAN} and \texttt{AsyncWrite} for buffering and compute/I/O
  overlap under BTD imbalance.
\item[$C_5$] Practical ADIOS2 configuration guidelines for WarpX-like PIC codes.
\end{description}

\subsection{Computational Artifacts}

\begin{description}
\item[$A_1$] \url{https://github.com/guj/artifacts26} ---
  Frontier and Perlmutter build scripts, WarpX inputs, PerformanceRun and
  \texttt{repeat/} jobs, log timer extractor, and result CSVs.
  (Optional Zenodo DOI may be added after acceptance.)
\end{description}

\begin{center}
\small
\begin{tabular}{@{}llp{0.48\columnwidth}@{}}
\toprule
Art. & Contrib. & Related paper elements \\
\midrule
$A_1$ & $C_1$--$C_5$ &
  Table~I; Fig.~2 (Perlmutter 3D); Figs.~3--6 (BTD);
  small-$N$ mean$\pm$stdev (\S II-B, \S III-C, \S IV-B);
  \S VI-B recommendations \\
\bottomrule
\end{tabular}
\end{center}

%%%%%%%%%%%%%%%%%%%%%%%%%%%%%%%%%%%%%%%%%%%%%%%%%%%%%%%%
\section{Artifact Identification}
%%%%%%%%%%%%%%%%%%%%%%%%%%%%%%%%%%%%%%%%%%%%%%%%%%%%%%%%

\newartifact

\artrel
$A_1$ packages inputs, ADIOS2 parameterizations, \texttt{repeat/} scripts,
the log timer extractor, and CSVs for Table~I and Figs.~2--6.
It supports replotting from CSVs and, where allocations permit, re-execution
of the WarpX~$\rightarrow$~openPMD-api~$\rightarrow$~ADIOS2 write path.

\artexp
\begin{itemize}
  \item Table~I (Frontier 3D; EWS/TLS $\pm$ Blosc2): relative ordering is the
        primary claim; absolute TB/s varies with Lustre contention.
  \item Fig.~2 (Perlmutter 3D): TLS above EWS; aggregate $>$1~TB/s at 512 nodes.
  \item BTD: \texttt{FlattenSteps} overhead negligible vs.\ EWS;
        \texttt{SPAN}+\texttt{AsyncWrite} highest among tested configs;
        \texttt{DataSizeBased} preferred when subfile count is constrained under skew.
\end{itemize}

\arttime
\begin{itemize}
  \item \textbf{Setup:} 30--120~min (build stack or load site modules).
  \item \textbf{Execution:} $\approx$1--2~min compute at 8 nodes (Slurm 5--10~min);
        queue wait is site-dependent. Full scale (2048 Frontier / 512 Perlmutter
        nodes) needs a leadership allocation and terabytes of scratch
        (Fig.~2: a 128-node Perlmutter 3D step is already $\approx$10~TB).
  \item \textbf{Analysis:} $<$15~min to import CSVs and replot in Google Sheets.
\end{itemize}

\artin

\artinpart{Hardware}
\begin{itemize}
  \item Frontier (ORNL; primary documented platform): HPE Cray EX; AMD MI250X;
        Lustre/ClusterStor ORION ($\approx$5~TB/s peak write); weak scaling to
        2048 nodes; 8 CPU cores/node for I/O. Setup and job scripts in the repo
        target Frontier.
  \item Perlmutter (NERSC): GPU nodes; all-flash Lustre ($>$5~TB/s); weak scaling
        to 512 nodes (scratch-quota limited); 4 CPU cores/node for I/O.
        Same stack via \path{scripts/perlmutter/build.setup} (CUDA).
        Example job: \path{warpx_tests/PerformanceRun/BTD/N8/perlmutter.sh}.
\end{itemize}
Node-local NVMe burst buffers were explored for BTD staging but are not required
to regenerate the primary CSV/plot artifacts.

\artinpart{Software}
\begin{itemize}
  \item WarpX 25.12 (\url{https://github.com/BLAST-WarpX/warpx})
  \item openPMD-api bundled with that WarpX build
        (\url{https://github.com/openPMD/openPMD-api})
  \item ADIOS2 2.11, BPFile / BP5 (\url{https://github.com/ornladios/ADIOS2})
  \item Frontier modules (\path{scripts/frontier/bash.setup.env}):
        \texttt{cmake/3.30.5}, \texttt{craype-accel-amd-gfx90a}, \texttt{rocm},
        \texttt{cray-mpich/8.1.31}, \texttt{cce/18.0.1}, \texttt{hdf5/1.14.5-mpi}
        (optional \texttt{ccache}, \texttt{ninja})
  \item Perlmutter modules (\path{scripts/perlmutter/build.setup}):
        \texttt{PrgEnv-gnu}, \texttt{craype-accel-nvidia80}, \texttt{cudatoolkit},
        \texttt{gcc-native/13.2}, \texttt{cmake/3.30.2}
  \item Lossless Blosc2 via ADIOS2 (\texttt{zstd} level~1, bit-shuffle,
        2048-byte threshold, 6 threads/rank)
  \item Python~3.6+ for \path{scripts/frontier/extract_time_3.6.py};
        Google Sheets to import \path{results/*.csv} (no matplotlib drivers)
\end{itemize}

\artinpart{Datasets / Inputs}
No external observational datasets. Decks and Slurm scripts live in
\texttt{warpx\_tests/}; three-trial scripts in \texttt{repeat/}.
Per-scale decks: \texttt{input.n\$\{NNODES\}}.
Scripts under \texttt{repeat/} resolve \texttt{EXE}, \texttt{bpls}, and
\texttt{inputs\_3d} relative to \texttt{warpx\_tests/} (set
\texttt{\#SBATCH -A YOUR\_PROJECT}).
\begin{itemize}
  \item \path{inputs_3d/opmd/regular/}, \texttt{regular\_blosc/},
        \texttt{regular\_h5/}: 3D Cartesian LWFA
        ($\approx$33M cells/node; $\approx$2.5~GB/node/step).
  \item \path{inputs_3d/opmd/btd_*}, \texttt{btd\_flatten*}: BTD
        (skew + out-of-order; FlattenSteps variants).
  \item \path{PerformanceRun/3D/}, \path{PerformanceRun/BTD/}: Frontier
        jobs (plus \path{BTD/.../perlmutter.sh}).
  \item \texttt{repeat/}: same pattern, by config
        (\texttt{bp/}, \texttt{bp\_async/}, \texttt{bp\_flatten/}, \texttt{h5/});
        submit each script three times for $N{=}8$--$32$
        (\path{repeat/README.md}).
\end{itemize}

\noindent Paper numbers are in \texttt{results/}:
{\small
\begin{center}
\begin{tabular}{@{}ll@{}}
\toprule
File & Paper \\
\midrule
\path{frontier_3d_throughput.csv} & Table~I \\
\path{3D-perlmutter.csv} & Fig.~2 \\
\path{fig_btd_a_flatten_frontier.csv} & Fig.~3 \\
\path{fig_btd_b_span_async_frontier.csv} & Fig.~4 \\
\path{fig_btd_c_agg_perlmutter.csv} & Fig.~5 \\
\path{fig_btd_d_subfiles_perlmutter.csv} & Fig.~6 \\
\path{repeat_runs_raw.csv}, \path{repeat_summary.csv} & mean$\pm$stdev \\
\bottomrule
\end{tabular}
\end{center}
}
See \path{warpx_tests_INDEX.md}, \path{repeat/README.md}, \path{results/README.md}.

\artinpart{Installation and Deployment}
Same tags on both sites (ADIOS2~2.11, WarpX~25.12, bundled openPMD-api,
\texttt{WarpX\_openpmd\_internal=ON}; override with \texttt{ADIOS2\_TAG},
\texttt{WARPX\_TAG}). Choose scratch \texttt{WORKDIR} (no hard-coded user path).
\begin{itemize}
  \item \textbf{Frontier (HIP):}\\
        \texttt{export WORKDIR=}\\
        \path{/lustre/orion/<project>/scratch/<user>/<exp>}\\
        \texttt{source scripts/frontier/bash.setup.env}\\
        \texttt{./scripts/frontier/build.setup}\\
        then e.g.\ \texttt{sbatch} \path{PerformanceRun/3D/N1/frontier_blosc.sh}
        or \texttt{sbatch} \path{PerformanceRun/BTD/N8/frontier.sh}.
  \item \textbf{Perlmutter (CUDA):}\\
        \texttt{export WORKDIR=}\\
        \path{/pscratch/sd/<user>/<exp>}\\
        \texttt{./scripts/perlmutter/build.setup} (loads NERSC modules)\\
        then e.g.\ \texttt{sbatch} \path{PerformanceRun/BTD/N8/perlmutter.sh}.
\end{itemize}
Then, on either host:
\begin{enumerate}
  \item In \texttt{warpx\_tests/}, symlink \texttt{EXE} $\rightarrow$ WarpX and
        \texttt{bpls} $\rightarrow$ ADIOS2 \texttt{bpls}; set \texttt{\#SBATCH -A}.
        Small-$N$ repeats: from \texttt{repeat/...}, submit each script three times.
  \item From the parent of each \texttt{job\_id/} directory:
        \path{python3 scripts/frontier/extract_time_3.6.py job_id}
        (bare id, no slash) $\rightarrow$ TotalTime, IOTime.
        Option~1 I/O time is
        $\mathrm{TotalTime}_{\mathrm{I/O}}-\mathrm{TotalTime}_{\mathrm{nullcore}}$
        when nullcore runs exist; \textbf{DataSize} is the observed on-disk
        output size.
  \item Aggregate into \path{results/*.csv}; replot in Google Sheets.
\end{enumerate}

\artcomp
$T_1$ build/load $\rightarrow$ $T_2$ run WarpX $\rightarrow$ $T_3$
\path{extract_time_3.6.py} plus file sizes into CSV $\rightarrow$ $T_4$
Google Sheets. Parameters: node count, aggregation, \texttt{NumSubFiles},
FlattenSteps, SPAN, AsyncWrite, Blosc2. Frontier $N{=}8$--$32$ configs were
repeated three times (\texttt{repeat/}); larger scales and Blosc are
single-trial, with reliability via cross-scale consistency (\S III-C).

\artout
CSV throughput vs.\ nodes plus small-$N$ mean$\pm$stdev; plots of Figs.~2--6
from those CSVs. Compare to Table~I and Figs.~2--6 on relative trends;
document filesystem contention as in \S VI-C.

\end{document}
